% Source PDF: source/普通高考/2007/2007湖北文.pdf, page 1, question 6.
% The source page is vector artwork (pdfimages -list found no embedded bitmap).
% Comparison source was rendered from the PDF after 100px/50px grid location.
% Grid evidence: tmp/redraw_grid/2007_hubei_liberal/page1_grid100.png and
% tmp/redraw_crops/2007_hubei_liberal/q6_block_grid50.png.
% Original comparison crop: tmp/redraw_crops/2007_hubei_liberal/q6_orig_chart.png.
% Elements: frequency-density axes with a broken x-axis, eleven width-2 kg
% histogram bars, twelve group-bound labels 54.5--76.5, and axis titles.
% Relations: each bar spans one 2 kg class interval; heights are the source
% frequency densities (0.01, 0.02, 0.05, 0.06, 0.07, 0.08, 0.065, 0.055,
% 0.04, 0.035, 0.015).  The displayed coordinates map one omitted group
% width to x=0..1, then map the twelve source group boundaries to x=1..12.
% solid segments: x/y axes, bar outlines, tick marks, and the x-axis break.
% dashed segments: none.  Group-bound labels remain horizontal; the chart is
% widened so the labels retain normal size without overlap.

\begin{tikzpicture}[baseline]
\begin{axis}[exam statistical histogram,scaled ticks=false,
    width=16.0cm,
    height=11.665cm,
    xmin=-0.95,
    xmax=14.0,
    ymin=-0.016,
    ymax=0.093,
    axis lines=none,
    ticks=none,
    clip=false,
    x tick label as interval=false,
]
    % Frequency-density bars; the final zero closes the last interval.
    \addplot[
        ybar interval,
        draw=black,
        fill=white,
        line width=0.65pt,
    ] coordinates {
        (1,0.01)
        (2,0.02)
        (3,0.05)
        (4,0.06)
        (5,0.07)
        (6,0.08)
        (7,0.065)
        (8,0.055)
        (9,0.04)
        (10,0.035)
        (11,0.015)
        (12,0)
    };

    % Coordinate axes.  The omitted 0--54.5 kg range is one displayed group width.
    \draw[->,line width=0.7pt] (axis cs:0,0) -- (axis cs:13.45,0);
    \draw[->,line width=0.7pt] (axis cs:0,0) -- (axis cs:0,0.090);

    % A pale, short zig-zag marks the broken part of the x-axis.
    \examstatxbreak[black!55,line width=0.55pt]{0.42000000000000004}{0};

    % Upright axis labels.
    \examstatylabel{\(\dfrac{\text{频率}}{\text{组距}}\)};
    \examstatxlabel{体重（kg）};
    \node[anchor=north east,inner sep=0pt,font=\normalsize,xshift=-4pt,yshift=-4pt] at (axis cs:-0.10,0.000000) {0};

    % Group-bound ticks and labels.  Each label is directly below its bar edge.
    \draw[line width=0.4pt] (axis cs:1,0) -- (axis cs:1,-0.0022);
    \node[anchor=north,inner sep=0pt,font=\normalsize,yshift=-2pt] at (axis cs:1,-0.002200) {54.5};
    \draw[line width=0.4pt] (axis cs:2,0) -- (axis cs:2,-0.0022);
    \node[anchor=north,inner sep=0pt,font=\normalsize,yshift=-2pt] at (axis cs:2,-0.002200) {56.5};
    \draw[line width=0.4pt] (axis cs:3,0) -- (axis cs:3,-0.0022);
    \node[anchor=north,inner sep=0pt,font=\normalsize,yshift=-2pt] at (axis cs:3,-0.002200) {58.5};
    \draw[line width=0.4pt] (axis cs:4,0) -- (axis cs:4,-0.0022);
    \node[anchor=north,inner sep=0pt,font=\normalsize,yshift=-2pt] at (axis cs:4,-0.002200) {60.5};
    \draw[line width=0.4pt] (axis cs:5,0) -- (axis cs:5,-0.0022);
    \node[anchor=north,inner sep=0pt,font=\normalsize,yshift=-2pt] at (axis cs:5,-0.002200) {62.5};
    \draw[line width=0.4pt] (axis cs:6,0) -- (axis cs:6,-0.0022);
    \node[anchor=north,inner sep=0pt,font=\normalsize,yshift=-2pt] at (axis cs:6,-0.002200) {64.5};
    \draw[line width=0.4pt] (axis cs:7,0) -- (axis cs:7,-0.0022);
    \node[anchor=north,inner sep=0pt,font=\normalsize,yshift=-2pt] at (axis cs:7,-0.002200) {66.5};
    \draw[line width=0.4pt] (axis cs:8,0) -- (axis cs:8,-0.0022);
    \node[anchor=north,inner sep=0pt,font=\normalsize,yshift=-2pt] at (axis cs:8,-0.002200) {68.5};
    \draw[line width=0.4pt] (axis cs:9,0) -- (axis cs:9,-0.0022);
    \node[anchor=north,inner sep=0pt,font=\normalsize,yshift=-2pt] at (axis cs:9,-0.002200) {70.5};
    \draw[line width=0.4pt] (axis cs:10,0) -- (axis cs:10,-0.0022);
    \node[anchor=north,inner sep=0pt,font=\normalsize,yshift=-2pt] at (axis cs:10,-0.002200) {72.5};
    \draw[line width=0.4pt] (axis cs:11,0) -- (axis cs:11,-0.0022);
    \node[anchor=north,inner sep=0pt,font=\normalsize,yshift=-2pt] at (axis cs:11,-0.002200) {74.5};
    \draw[line width=0.4pt] (axis cs:12,0) -- (axis cs:12,-0.0022);
    \node[anchor=north,inner sep=0pt,font=\normalsize,yshift=-2pt] at (axis cs:12,-0.002200) {76.5};

    % Frequency-density ticks and labels.
    \draw[line width=0.4pt] (axis cs:0,0.01) -- (axis cs:-0.14,0.01);
    \node[anchor=east,inner sep=0pt,font=\normalsize,xshift=-4pt] at (axis cs:-0.140000,0.01) {0.01};
    \draw[line width=0.4pt] (axis cs:0,0.02) -- (axis cs:-0.14,0.02);
    \node[anchor=east,inner sep=0pt,font=\normalsize,xshift=-4pt] at (axis cs:-0.140000,0.02) {0.02};
    \draw[line width=0.4pt] (axis cs:0,0.03) -- (axis cs:-0.14,0.03);
    \node[anchor=east,inner sep=0pt,font=\normalsize,xshift=-4pt] at (axis cs:-0.140000,0.03) {0.03};
    \draw[line width=0.4pt] (axis cs:0,0.04) -- (axis cs:-0.14,0.04);
    \node[anchor=east,inner sep=0pt,font=\normalsize,xshift=-4pt] at (axis cs:-0.140000,0.04) {0.04};
    \draw[line width=0.4pt] (axis cs:0,0.05) -- (axis cs:-0.14,0.05);
    \node[anchor=east,inner sep=0pt,font=\normalsize,xshift=-4pt] at (axis cs:-0.140000,0.05) {0.05};
    \draw[line width=0.4pt] (axis cs:0,0.06) -- (axis cs:-0.14,0.06);
    \node[anchor=east,inner sep=0pt,font=\normalsize,xshift=-4pt] at (axis cs:-0.140000,0.06) {0.06};
    \draw[line width=0.4pt] (axis cs:0,0.07) -- (axis cs:-0.14,0.07);
    \node[anchor=east,inner sep=0pt,font=\normalsize,xshift=-4pt] at (axis cs:-0.140000,0.07) {0.07};
    \draw[line width=0.4pt] (axis cs:0,0.08) -- (axis cs:-0.14,0.08);
    \node[anchor=east,inner sep=0pt,font=\normalsize,xshift=-4pt] at (axis cs:-0.140000,0.08) {0.08};

\end{axis}
\end{tikzpicture}
